% !TeX root = main.tex
\section{策略评价与应用}\label{sec:evaluation}

给定偏好下的策略评价要回到最初定义的长期决策准则，训练损失、内部奖励和解集指标反映的是算法的部分性质，最终要判断的是策略是否实现了预期的目标取舍和约束要求。本章的安排是：先讨论单个策略与策略集合在原准则下的评价；随后考察评价所依赖的信息条件，包括只有历史数据时的离线评价、实验比较规范以及鲁棒性与部署条件；最后结合代表性应用，说明不同层次的实验证据各自能支持哪些结论。

\subsection{单策略在原准则下的评价}\label{sec:eval-criteria}

对于采用实值效用的任务，首先应报告策略在原长期准则下的表现，同时保留各目标的结果，以便分析效用的改善来自哪些目标、为此付出了什么代价。能精确求得最优值时，可以进一步报告最优性差距；只有基线或近似参考解时，结果应表述为相对于该参考的性能差异。严格词典序等非标量比较关系，按原比较规则评价，不额外构造加权总分。

当评价对象为累计回报向量时，SER 与 ESR 对样本的聚合顺序不同。设同一固定策略在相同初始状态分布下独立执行，得到累计回报样本 \(\mathbf G^{(1)},\ldots,\mathbf G^{(N)}\)，两者的估计为
\begin{equation}
\begin{aligned}
\widehat J_{\mathrm{SER}}
&=
u_p\!\left(
\frac{1}{N}
\sum_{n=1}^{N}\mathbf G^{(n)}
\right),\\
\widehat J_{\mathrm{ESR}}
&=
\frac{1}{N}
\sum_{n=1}^{N}
u_p\!\left(\mathbf G^{(n)}\right).
\end{aligned}
\end{equation}
SER 先对回报向量取样本均值再施加效用，ESR 先计算每次执行的效用再取平均，与第~\ref{sec:criteria-stochastic}~节的定义的运算顺序一致\cite{hayes2022}。非线性效用下，SER 的有限样本代入估计有其自身的统计性质，置信区间和误差分析也要按相同的聚合顺序进行。

上式对应以累计回报向量为评价对象的准则。若原目标是阶段效用的期望累积，就先在每条轨迹内按既定时域和折扣累加阶段效用，再对轨迹结果取期望；若准则直接作用于轨迹特征，就先从完整轨迹或相应的充分统计量中计算该特征，再进行评价。路径事件、首次到达时间和累计资源使用等指标同样按各自的定义处理。评价公式由原长期准则决定，而不是由算法的输出形式决定。

满意水平、风险和约束问题，除效用本身外还要报告达标与违反情况。多个目标需要同时满足时，统计联合达标概率；连续约束可以同时报告违反频率和违反幅度。罕见失效或尾部风险的结果，要结合样本规模、评价时域和不确定性估计来解释。

除数值比较外，还可以展示策略的目标分量、主要取舍和结果的不确定性。多目标决策支持的研究已经系统整理了候选解可视化、解集分析和不确定性探索方法\cite{osika2023decisionsupport}；用于序贯策略时，这些展示的作用是解释给定偏好下的选择依据，并与原准则下的实测表现对应起来。

\subsection{策略集合与条件策略的评价}\label{sec:eval-sets}

算法输出策略库、近似 Pareto 前沿或偏好条件策略时，解集质量与最终策略质量需要分开评价。解集指标描述候选集合的覆盖、分布或逼近质量。Zitzler 等系统分析了近似解集质量指标的表达能力与局限\cite{zitzler2003}，这些指标也广泛用于多目标强化学习的评价\cite{vamplew2011}。超体积刻画了解集在目标空间中的覆盖程度，但不说明某个给定偏好下所选策略的效用。

解集指标本身也可能带有隐含的取舍。加权超体积把候选集合支配区域内的体积积分改为 $\int_{D(\mathcal L,\mathbf z_{\rm ref})}h(\mathbf z)\,d\mathbf z$，其中 $h$ 是目标空间上的权重密度，用来改变不同结果区域对指标的贡献\cite{zitzler2007weighted}。这种构造可以突出极端结果、参考区域或特定目标，密度、目标尺度和参考点的选择则需要给出依据。目标空间上的积分权重与偏好空间上的用户分布是两个不同的对象，加权超体积也不等于某个给定效用下的策略表现。

对于采用实值评价 $J_p$ 且候选满足可接受性条件的任务，策略库 $\mathcal L$ 在偏好分布 $q$ 下的集合期望效用，评价的是从库中选取最好候选能达到的表现；偏好条件策略的评价则应使用输入 $p$ 后实际输出的策略 $\pi_p$。二者分别为
\begin{equation}
\operatorname{EU}_q(\mathcal L)
=\mathbb E_{p\sim q}\!\left[\max_{\pi\in\mathcal L}J_p(\pi)\right],
\qquad
\mathbb E_{p\sim q}[J_p(\pi_p)].
\end{equation}
即使输出的候选集合不变，交换各策略对应的偏好标签也可能保持超体积和集合期望效用不变，却改变实际输入偏好所得到的效用。因此，解集评价要和偏好条件下的最终策略评价一起报告，并说明偏好的采样分布以及局部区域的性能差异。Felten 等提供的期望效用和参考效用损失指标，为这种评价提供了可用的框架\cite{felten2023}。


期望效用依赖偏好分布，最大效用损失则考察指定偏好族中最不利的近似误差。Zintgraf 等以用户效用为依据提出这两类评价，并将几何覆盖误差与效用损失联系起来\cite{zintgraf2015quality}。对于参考策略集 $\mathcal R$，相应的损失可以写为
\begin{equation}
\operatorname{MUL}_{\mathcal P}(\mathcal L;\mathcal R)
=\sup_{p\in\mathcal P}\left[
\max_{\pi\in\mathcal R}J_p(\pi)
-\max_{\pi\in\mathcal L}J_p(\pi)\right].
\end{equation}
若评价对象为期望回报向量，效用对各目标单调且在欧氏范数下具有统一的 Lipschitz 常数 $L$，候选集对真实 Pareto 前沿达到逐坐标加性 $\epsilon$ 覆盖，则最大效用损失不超过 $L\sqrt m\,\epsilon$。这一结论属于 SER 语义，不适用于仅由均值构造的 ESR 评价。参考集若只是多个算法经验结果的并集，所得数值是相对于该参考的效用差异，而不是真实的最优性差距。平均效用、最不利偏好损失与几何覆盖互为补充，报告时要一并说明偏好族、尺度和参考信息。

偏好输入与实际输出之间的响应关系，可以作为补充诊断。de las Heras Molins 等以每个目标的权重与对应回报之间的 Spearman 秩相关刻画条件策略的可控性，检验集合指标没有记录的偏好—结果配对关系\cite{heras2026controllability}。这一指标假定相应目标随权重增加而改善，也需要规定其他权重如何随之调整。相关性高说明可控性好，但不能代替原效用、最优性差距或约束满足的评价；多个偏好共享同一个最优策略时，局部输出不变也可能完全符合原准则。

偏好条件策略还要区分偏好空间内的插值与训练范围之外的外推。训练时覆盖的偏好区域决定了模型能直接获得哪些监督或交互信息，部署偏好离这一范围越远，策略越依赖函数逼近的泛化能力。Zhu 等在不同偏好覆盖条件下的实验表明，偏好分布本身就是多目标离线学习的重要数据条件\cite{zhu2023}。评价条件策略时，可以分别报告训练偏好范围内、边界区域和未见偏好下的效用与约束表现，并与针对单一偏好专门训练的策略比较。

训练标签的组织方式也会影响偏好泛化。针对第~\ref{sec:solving-reuse}~节讨论的回溯重标记，Bonin 等以相对基线的效用损失超过规定阈值的偏好比例，衡量偏好被“放弃”的程度\cite{bonin2026collapse}。这一诊断依赖参考基线、损失阈值和偏好采样分布，能揭示平均指标没有反映出来的局部退化；其证据来自在线采样与经验回放，固定离线日志下的情形还需单独检验。

\subsection{离线策略评价}\label{sec:eval-offline}

可用于评价的信息随数据条件变化。有独立仿真器或受控实验环境时，可以直接执行策略获得评价样本，但要说明评价环境与实际部署条件的差异。只有历史数据时，要采用离线策略评价方法，并写明其覆盖和估计条件。

用同一批轨迹比较不同目标与策略，已有较早的研究基础。Shelton 利用策略之间的似然比估计多目标回报曲面，结合带记忆的策略表示与参数化策略搜索，分析了多个期望结果之间的权衡及约束优化\cite{shelton2001importance}。这种复用依赖轨迹概率与回报的记录，可靠的估计区域受行为数据支持和方差的限制。可见，多目标离线评价既包括更换目标权重，也包括检验候选策略是否得到已有数据的支持。

典型的离线策略评价方法有重要性加权、模型估计和价值函数估计。重要性加权要求行为策略覆盖目标策略的动作，模型与价值估计则要处理拟合误差。Jiang 和 Li 的双重稳健估计把价值预测与重要性加权残差校正结合起来\cite{jiang2016}。对于长度为 $T$ 的轨迹，逐目标应用其递归形式可写为
\begin{equation}
D_{t,i}
=\widehat V_{t,i}(s_t)
+\rho_t\left[r_{t,i}+\gamma D_{t+1,i}
-\widehat Q_{t,i}(s_t,a_t)\right],\qquad D_{T,i}=0,
\end{equation}
其中 $\rho_t=\pi_t(a_t\mid s_t)/b_t(a_t\mid s_t)$，$b$ 为行为策略，$\widehat V_{t,i}(s)=\mathbb E_{a\sim\pi_t}[\widehat Q_{t,i}(s,a)]$。式中的状态应包含策略所需的信息：对于依赖累计回报或历史的策略，概率比要在相应的增广状态或历史上计算，日志也必须能恢复这些条件变量。目标策略和拟合的价值函数应与用于检验的轨迹分离。在行为概率已知、目标策略有支持等条件下，准确的重要性比率使该估计即使在价值预测不准时也保持均值无偏；价值预测的作用是减小残差波动，但并非任何拟合器都能降低方差。逐目标估计得到的是期望向量，它既不恢复目标之间的联合分布，代入非线性效用后得到的 SER 估计也不再无偏。

ESR、分位数、CVaR 和轨迹约束需要的不只是期望回报。UnO 从标量回报的累积分布函数出发，其基本估计为
\begin{equation}
\widehat F(z)=\frac{1}{N}\sum_{n=1}^{N}
\rho^{(n)}\,\mathbf 1\{G^{(n)}\leq z\},
\end{equation}
其中 $\rho^{(n)}$ 是整条轨迹上目标策略与行为策略的概率比。方法先构造分布函数的置信带，再在与该置信带相容的分布上求所需统计量的上下界，从而统一处理均值、分位数和 CVaR 等评价\cite{chandak2021}。其信息条件包括独立轨迹、已知的行为动作概率、目标策略有支持以及回报有界；长时域的概率比可能带来较大的方差。单个标量回报的分布不包含多个目标之间的相依关系；对于已知的固定 ESR，若日志足以重建每条轨迹的效用，可以把该效用当作标量评价对象。

标量分布评价还可以借助模型预测降低重要性加权的方差。OPRA 对有限时域 MDP 的回报分布建立递归估计，把模型预测的条件分布函数作为基线，再用实际奖励和后继样本做加权残差校正\cite{huang2022opra}。在回报有界、行为概率准确、目标策略有支持，且条件分布模型已固定或独立于评价轨迹拟合等条件下，原始的双重稳健估计保持分布函数的逐点无偏；方差能否降低，取决于条件分布模型的误差。原始估计可能不单调或超出 $[0,1]$，实际使用时需要单调化和裁剪。对于在分布函数上满足 Lipschitz 条件的风险泛函 $\rho$，一致的分布误差可以转化为
\begin{equation}
|\rho(\widehat F)-\rho(F)|
\leq L_\rho\|\widehat F-F\|_\infty,
\end{equation}
这为同一分布估计下的多种风险评价提供了共同的误差依据，例如 CVaR 的相应常数还取决于回报范围与尾部概率。原始分布估计无偏，并不意味着非线性风险的代入估计无偏；标量结果也不提供多个目标联合事件的保证。

需要保留多个目标的联合长期结果时，可以直接做多变量分布评价。Qi 等为给定目标策略建立条件联合回报分布的 Bellman 方程，用接受状态、动作和噪声输入的生成器逼近该分布\cite{qi2025distributional}。训练时用 Wasserstein 距离比较生成分布与由当前奖励、下一状态及目标策略形成的分布备份，并利用分布 Bellman 算子的收缩性分析估计误差。数据由固定的平稳行为策略产生，被评价的目标策略也是平稳策略。有限样本结果依赖时齐 Markov 结构、已知有界的奖励机制，以及给定状态动作后奖励与下一状态条件独立；同时要求状态动作访问密度有正下界、几何遍历性，以及生成与判别函数满足相应的光滑性和网络复杂度条件。分布误差可以控制有界 Lipschitz 效用的期望误差，不连续的阈值事件、分位数或尾部风险则没有同等的置信保证。

三类方法提供的信息不同：标量均值估计给出期望向量，标量分布置信带给出单个评价量的分布，多变量联合分布估计保留了目标之间的相依关系。选择哪一类，由原长期准则需要什么信息决定。

表~\ref{tab:evaluation-information}~进一步列出原准则所需的评价对象、可用路线与离线条件。
\begin{table}[H]
\centering
\caption{原长期准则所需的评价信息与离线估计条件}
\label{tab:evaluation-information}
\begingroup
\zihao{5}\songti
\setlength{\tabcolsep}{5pt}
\renewcommand{\arraystretch}{1.18}
\begin{tabularx}{\textwidth}{@{}P{0.15\textwidth}Y Y Y@{}}
\toprule
准则 & 评价对象 & 估计路线 & 关键边界 \\
\midrule
SER & 期望累计回报向量 & 逐目标估计均值，再代入效用\cite{jiang2016} & 非线性代入不自动无偏 \\
\addlinespace[6pt]
固定效用 ESR & 轨迹效用或充分信息 & 直接评价轨迹效用；必要时使用分布信息 & 均值不足；不必完整联合分布 \\
\addlinespace[6pt]
联合达标 & 同次轨迹的联合事件 & 事件概率，或适配的联合结果模型\cite{qi2025distributional} & 边缘概率不足；分布误差不自动保证阈值事件 \\
\addlinespace[6pt]
CVaR 等\newline 尾部指标 & 指定标量损失的分布与尾部 & 分布置信带，或分布估计后计算风险\cite{chandak2021,huang2022opra} & 明确损失、尾部概率、有界性和风险泛函条件 \\
\addlinespace[6pt]
轨迹特征\newline 或路径要求 & 事件、峰值、时长等统计量 & 从日志恢复轨迹特征，按原准则评价 & 累计向量未必保留发生次序与持续时长 \\
\bottomrule
\end{tabularx}
\par\vspace{4pt}\raggedright\noindent
注：共同检查目标策略支持、日志记录、尺度与初始分布、模型或函数类假设，以及训练和评价数据的关系。轨迹效用或事件的标量评价须匹配轨迹或扩展状态表示；历史依赖策略的概率比也须以相应信息为条件。
\par\endgroup
\end{table}


\subsection{实验比较与统计可靠性}\label{sec:eval-stats}

实验比较要统一问题定义和评价设置，包括目标方向、归一化方式、时域、折扣因子、初始状态分布和约束定义，其中任何一项变化都可能改变策略排序。MO-Gymnasium、MORL-Baselines 等统一接口减少了实现差异，为多目标强化学习实验提供了共享基准\cite{felten2023}；但共享环境并不保证各方法采用相同的效用语义。特别是使用线性权重时，如果不同方法对目标尺度的处理不同，相同的数值权重会对应不同的实际取舍。

环境随机性应作为明确的实验维度。Vamplew 等在若干简单多目标 MDP 中分别考察了随机奖励、随机转移以及 SER/ESR 目标对价值学习表现的影响，表明只在确定性任务上比较会漏掉重要差异\cite{vamplew2022stochasticity}。实验分析时要区分两种情况：随机性改变了策略的结果分布，还是暴露了价值表示与原准则之间的不匹配。确定性基准上的解集覆盖，不代表随机环境中的偏好一致性。

资源条件也应纳入比较。模型查询、真实环境交互、离线数据规模、预训练和超参数搜索都是求解成本。跨偏好方法适合同时报告预训练总成本、获得一个新偏好策略的追加成本和部署阶段的计算开销；离线方法则要区分用固定数据训练、用验证数据选择模型和额外调用环境信息调参这三种信息使用。

随机性来自算法训练和策略执行两个层次。对单个训练模型做大量测试，提高的是该模型性能估计的精度，反映不了不同随机种子下训练结果的稳定性\cite{vamplew2011,agarwal2021}。因此，应分别报告独立训练次数、每个训练实例的测试轨迹数和最终的汇总方式。数据来自多个工况、设备或时间段时，应按合适的独立单位划分训练和测试，避免高度相关的样本造成过于乐观的估计。

基线和消融实验要与研究问题对应。能直接线性标量化的问题，标量化后的单目标方法是必要的基线；新方法若依赖扩展状态、偏好条件输入、分布价值或特殊约束机制，可以通过消融分析这些设计对原准则的贡献。工程应用中，还应保留有实际意义的规则控制、传统优化方法或现有控制器作为参照。报告性能差异时，同时说明数据、资源与问题设置，并用相应的实验检验对差异成因的解释；多次独立实验和完整的结果分布，比单次最优结果更能反映算法的实际表现\cite{agarwal2021}。

\subsection{鲁棒性与部署条件下的表现}\label{sec:eval-robust}

训练分布上的性能只反映策略在既定条件下的表现。实际部署中，动力学参数、观测噪声、需求模式和初始状态都可能变化，因此还要考察这些变化对效用和约束满足的影响。

鲁棒性首先要明确不确定性的来源和评价方式。Ide 和 Schöbel 比较了多目标优化中的多种鲁棒有效性概念，表明不同的不确定性模型和比较规则对应不同的鲁棒性含义\cite{ide2016}。在序贯决策中，还要区分不确定性作用于状态转移、观测还是回报，以及采用期望、最坏情形、风险度量还是其他评价准则。扰动后的评价仍沿用原长期准则，而不是换成某个便于测量的单一目标。

部署还涉及实时性、资源消耗和失效处理。规划方法要在决策周期内完成在线优化，学习策略则把主要计算放在训练阶段。比较运行时间时，应区分训练成本、单步推理成本和在线规划成本，并在相同硬件和任务规模下报告平均延迟及必要的高分位延迟。系统若包含安全过滤、约束投影等动作修正模块，最终评价的应是修正后实际执行的策略。

\subsection{代表性应用中的建模与验证}\label{sec:eval-apps}

应用研究把偏好、长期准则和求解机制落实到具体任务中，也让评价结论的适用范围更清楚。下面结合基础设施维护、节能驾驶、临床决策、水库运行和设备控制，比较非线性长期效用、偏好条件策略和历史数据学习在不同信息条件下的处理方式。

基础设施维护将非线性效用与部分观测结合：MO-DCMAC 以信念描述退化状态，以累计结果和多变量回报分布支持 ESR 决策\cite{remmerden2025}。其模型案例的成本与倒塌概率取舍，不能直接视为长期现场部署的证据。

节能驾驶中的 MOAC 以交通状态和偏好为条件生成策略，在 SUMO 中比较行驶效率与能耗\cite{he2023}。训练的偏好对齐项与碰撞惩罚须分别解释，后者属于软性处理，不能据此宣称满足硬安全约束。

CATIE 案例用多奖励线性 FQI 分析序贯治疗，展示不同症状、副作用与生活质量取舍下的治疗建议\cite{lizotte2012}。这是历史试验数据上的回顾性策略分析，推荐策略的前瞻性临床获益仍需另行验证。

水库运行中的 MOFQI 以权重为条件共享价值与放水规则\cite{castelletti2013water}。训练样本由模型结合历史入流生成，测试也在历史入流条件下仿真，因此不能将其信息条件写成仅有被动运行日志。

服务器功率控制用固定系数表达响应性能与功耗的取舍，并在刀片服务器集群上检验日志训练的策略\cite{tesauro2007power}。其目标是阶段效用的折扣累积，区别于累计结果上的非线性 ESR。

湿式离合器研究比较接合时间与接合时活塞速度的不同标量化及参数搜索\cite{brys2013clutch}。该篇明确将实验限制在仿真；其引用的早期工作具有实机验证，二者的证据不能合并。表~\ref{tab:application-evidence}~列出六个案例的建模与验证条件。

\begin{table}[H]
\centering
\caption{代表性应用的决策建模与验证条件}
\label{tab:application-evidence}
\begingroup
\zihao{5}\songti
\setlength{\tabcolsep}{5pt}
\renewcommand{\arraystretch}{1.18}
\setlength{\reviewtablewidth}{\dimexpr\textwidth-6\tabcolsep\relax}
\begin{tabular}{@{}P{0.14\reviewtablewidth}P{0.30\reviewtablewidth}P{0.30\reviewtablewidth}P{0.26\reviewtablewidth}@{}}
\toprule
应用 & 目标、准则与表示 & 数据与验证 & 证据边界 \\
\midrule
基础设施\newline 维护\cite{remmerden2025} & 成本与倒塌概率；非线性 ESR；信念、累计量与联合分布 & 三类退化模型；含 Amsterdam 码头墙简化模型 & 模型环境中的维护决策；非长期现场部署 \\
\addlinespace[7pt]
节能驾驶\cite{he2023} & 效率与能耗；线性权重条件策略；训练含偏好对齐项 & SUMO 高速公路仿真；多种交通密度 & 仿真内偏好响应；碰撞惩罚非硬安全约束 \\
\addlinespace[7pt]
序贯治疗\cite{lizotte2012} & 症状、BMI、生活质量；线性取舍与多奖励 FQI & CATIE 随机试验简化子集；两阶段回顾分析 & 历史数据下的策略比较；非前瞻性临床获益 \\
\addlinespace[7pt]
水库运行\cite{castelletti2013water} & 发电与防洪；固定线性权重；条件价值与放水规则 & 模型和 1957--1978 年入流训练；1997--2004 年入流仿真检验 & 真实水库背景的模型实验；非纯被动日志学习 \\
\addlinespace[7pt]
服务器\newline 功率控制\cite{tesauro2007power} & 响应时间与功耗；固定取舍；阶段效用折扣累积 & 日志训练；刀片服务器集群与规定负载试验 & 有设备执行证据；非累计结果的非线性 ESR \\
\addlinespace[7pt]
湿式\newline 离合器\cite{brys2013clutch} & 接合时间与活塞速度；三类标量化；开环参数搜索 & 填充阶段 Simulink 模型；本篇仅仿真 & 标量化与搜索的影响；前作实机验证非本篇证据 \\
\bottomrule
\end{tabular}
\par\vspace{4pt}\raggedright\noindent
注：比较问题建模与证据条件，不作跨任务排名。真实背景、历史数据、仿真、设备试验与长期部署分别记录；具体准则与算法条件见正文及原文。
\par\endgroup
\end{table}


这些案例的长期准则、数据来源与验证环境各不相同。评价应用研究时，需要同时检查偏好如何定义、算法实际优化什么，以及实验支持哪些结论；真实系统背景、历史数据、设备试验与长期部署应分别记录。同样的“成本”“风险”或“安全”可能对应不同的数学对象，比较前仍须统一决策语义。由此产生的表示复杂度、离线评价、分布变化和部署问题，将在下一章讨论。
